\documentclass[fleqn]{article}
\usepackage{amsmath}
\usepackage{amsfonts}
\usepackage{enumitem}
\title{Solución Problemas Números Complejos}
\author{Juan Rodríguez}

\setlength{\parindent}{0pt}

\begin{document}
\maketitle
\section{Ejercicio 1}
	Utilizando el Teorema Fundamental del Cálculo, determina las derivadas de las siguientes funciones:

	\[
	F_1(x)=\int_0^x t^2\,dt
	\qquad
	F_2(x)=\int_0^{x^2} e^t\,dt
	\]

	\[
	F_3(x)=\int_1^{e^{3x}}\sin(t)\,dt
	\qquad
	F_4(x)=\int_2^{\sin(x)}\frac{1}{1-t^2}\,dt
	\]

	\subsection{Solución}

	Aplicamos el Teorema Fundamental del Cálculo:

	\[
	F'(x)=f(g(x))g'(x)
	\]

	\subsubsection*{a)}

	\[
	F_1'(x)
	=
	x^2
	\]

	\[
	\boxed{F_1'(x)=x^2}
	\]

	\subsubsection*{b)}

	\[
	F_2'(x)
	=
	e^{x^2}(x^2)'
	=
	2xe^{x^2}
	\]

	\[
	\boxed{F_2'(x)=2xe^{x^2}}
	\]

	\subsubsection*{c)}

	\[
	F_3'(x)
	=
	\sin(e^{3x})(e^{3x})'
	=
	3e^{3x}\sin(e^{3x})
	\]

	\[
	\boxed{F_3'(x)=3e^{3x}\sin(e^{3x})}
	\]

	\subsubsection*{d)}
	\[
	F_4'(x)
	=
	\frac{1}{1-\sin^2(x)}\cos(x)
	=
	\frac{\cos(x)}{\cos^2(x)}
	=
	\frac{1}{\cos(x)}
	\]

	\[
	\boxed{F_4'(x)=\frac{1}{\cos(x)}}
	\]

\section{Ejercicio 2}
	Calcula la ecuación de la recta tangente en $x=1$ a la gráfica de

	\[
	F(x)=\int_{-1}^{x}\frac{t^3}{t^4-4}\,dt
	\]

	\subsection{Solución}

	La ecuación de la recta tangente es:

	\[
	y-F(1)=F'(1)(x-1)
	\]

	Por el Teorema Fundamental del Cálculo:

	\[
	F'(x)=\frac{x^3}{x^4-4}
	\]

	\[
	F'(1)=\frac{1}{1-4}=-\frac{1}{3}
	\]

	Calculamos $F(1)$:

	\[
	F(1)
	=
	\int_{-1}^{1}\frac{t^3}{t^4-4}\,dt
	=
	0
	\]

	ya que el integrando es una función impar.

	Por tanto:

	\[
	y
	=
	-\frac{1}{3}(x-1)
	\]

	\[
	\boxed{
	y=-\frac{x}{3}+\frac{1}{3}
	}
	\]

\section{Ejercicio 3}
	Determina los intervalos en los que la función

	\[
	F(x)=\int_1^x \arctan(e^t)\,dt
	\]

	es inyectiva.

	\subsection{Solución}

	Por el Teorema Fundamental del Cálculo:

	\[
	F'(x)=\arctan(e^x)
	\]

	Como:

	\[
	e^x>0
	\]

	para todo $x\in\mathbb{R}$, entonces:

	\[
	\arctan(e^x)>0
	\]

	Por tanto:

	\[
	F'(x)>0
	\qquad
	\forall x\in\mathbb{R}
	\]

	Luego $F(x)$ es estrictamente creciente en todo su dominio y, por tanto, es inyectiva en:

	\[
	\boxed{(-\infty,+\infty)}
	\]

\section{Ejercicio 4}
	Calcula

	\[
	\lim_{x\to 0}
	\frac{1}{x}
	\int_0^x
	\frac{|\cos(t^3)|}{t^2+1}\,dt
	\]

	\subsection{Solución}

	Definimos:

	\[
	F(x)
	=
	\int_0^x
	\frac{|\cos(t^3)|}{t^2+1}\,dt
	\]

	Como:

	\[
	F(0)=0
	\]

	el límite presenta una indeterminación $\frac{0}{0}$. Aplicamos la regla de L'Hôpital y el Teorema Fundamental del Cálculo:

	\[
	\lim_{x\to 0}
	\frac{F(x)}{x}
	=
	\lim_{x\to 0}
	\frac{F'(x)}{1}
	=
	\lim_{x\to 0}
	\frac{|\cos(x^3)|}{x^2+1}
	\]

	\[
	=
	\frac{|\cos(0)|}{1}
	=
	1
	\]

	\[
	\boxed{1}
	\]

\section{Ejercicio 5}
	Calcula, en caso de existir, la primera y segunda derivada de la función

	\[
	F(x)=x\int_{2x}^{3x}e^{-t^2}\,dt
	\]

	\subsection{Solución}

	Aplicamos la regla del producto:

	\[
	F'(x)
	=
	\int_{2x}^{3x}e^{-t^2}\,dt
	+
	x\frac{d}{dx}
	\left(
	\int_{2x}^{3x}e^{-t^2}\,dt
	\right)
	\]

	Utilizando el Teorema Fundamental del Cálculo:

	\[
	\frac{d}{dx}
	\left(
	\int_{2x}^{3x}e^{-t^2}\,dt
	\right)
	=
	3e^{-9x^2}
	-
	2e^{-4x^2}
	\]

	Por tanto:

	\[
	\boxed{
	F'(x)
	=
	\int_{2x}^{3x}e^{-t^2}\,dt
	+
	3xe^{-9x^2}
	-
	2xe^{-4x^2}
	}
	\]

	Derivamos de nuevo:

	\[
	F''(x)
	=
	3e^{-9x^2}
	-
	2e^{-4x^2}
	+
	3e^{-9x^2}
	-
	54x^2e^{-9x^2}
	-
	2e^{-4x^2}
	+
	16x^2e^{-4x^2}
	\]

	\[
	F''(x)
	=
	6e^{-9x^2}
	-
	4e^{-4x^2}
	-
	54x^2e^{-9x^2}
	+
	16x^2e^{-4x^2}
	\]

	\[
	\boxed{
	F''(x)
	=
	(6-54x^2)e^{-9x^2}
	+
	(16x^2-4)e^{-4x^2}
	}
	\]

\section{Ejercicio 6}
	Demuestra que la función

	\[
	F(x)=\int_{1-x}^{1+x}\ln(t)\,dt
	\]

	es decreciente en $\left(0,\frac{1}{2}\right)$.

	\subsection{Solución}

	Aplicamos el Teorema Fundamental del Cálculo:

	\[
	F'(x)
	=
	\ln(1+x)(1)
	-
	\ln(1-x)(-1)
	\]

	\[
	F'(x)
	=
	\ln(1+x)+\ln(1-x)
	=
	\ln\left((1+x)(1-x)\right)
	=
	\ln(1-x^2)
	\]

	Para $x\in\left(0,\frac{1}{2}\right)$:

	\[
	0<1-x^2<1
	\]

	y, por tanto:

	\[
	\ln(1-x^2)<0
	\]

	\[
	F'(x)<0
	\qquad
	\forall x\in\left(0,\frac{1}{2}\right)
	\]

	Por tanto:

	\[
	\boxed{
	F(x)\text{ es decreciente en }
	\left(0,\frac{1}{2}\right)
	}
	\]

\section{Ejercicio 7}
	Determina el polinomio de Taylor de orden 3 en $x_0=0$ de la función

	\[
	F(x)=\int_0^x t^2\cos(t^2)\,dt
	\]

	y utiliza el resultado para calcular

	\[
	\lim_{x\to0}\frac{F(x)}{x^3}
	\]

	\subsection{Solución}

	Por el Teorema Fundamental del Cálculo:

	\[
	F'(x)=x^2\cos(x^2)
	\]

	\[
	F''(x)
	=
	2x\cos(x^2)-2x^3\sin(x^2)
	\]

	\[
	F'''(x)
	=
	2\cos(x^2)
	-10x^2\sin(x^2)
	-4x^4\cos(x^2)
	\]

	Calculamos los valores en $x=0$:

	\[
	F(0)=0,
	\qquad
	F'(0)=0,
	\qquad
	F''(0)=0,
	\qquad
	F'''(0)=2
	\]

	Por tanto, el polinomio de Taylor de orden 3 es:

	\[
	P_3(x)
	=
	F(0)+F'(0)x+\frac{F''(0)}{2!}x^2
	+\frac{F'''(0)}{3!}x^3
	=
	\frac{x^3}{3}
	\]

	\[
	\boxed{P_3(x)=\frac{x^3}{3}}
	\]

	Entonces:

	\[
	F(x)
	=
	\frac{x^3}{3}+o(x^3)
	\]

	\[
	\lim_{x\to0}\frac{F(x)}{x^3}
	=
	\lim_{x\to0}
	\frac{\frac{x^3}{3}+o(x^3)}{x^3}
	=
	\frac{1}{3}
	\]

	\[
	\boxed{\frac{1}{3}}
	\]

\section{Ejercicio 8}
	Dada la función $f(x)=x^2$, determina para qué valor
	$c\in[0,3]$ se verifica el teorema del valor medio integral.

	\subsection{Solución}

	Por el teorema del valor medio integral:

	\[
	\int_a^b f(x)\,dx
	=
	f(c)(b-a)
	\]

	En este caso:

	\[
	\int_0^3 x^2\,dx
	=
	c^2(3-0)
	\]

	\[
	\left[\frac{x^3}{3}\right]_0^3
	=
	3c^2
	\]

	\[
	9=3c^2
	\]

	\[
	c^2=3
	\]

	Como $c\in[0,3]$:

	\[
	\boxed{c=\sqrt{3}}
	\]

\section{Ejercicio 9}
	Utiliza el teorema del valor medio del cálculo integral para determinar el valor
	(o valores) $c\in(a,b)$ que lo satisface(n) respecto a la función

	\[
	f(x)=1-x+x^2
	\]

	en el intervalo $[-1,1]$.

	\subsection{Solución}

	Por el teorema del valor medio integral:

	\[
	\int_a^b f(x)\,dx
	=
	f(c)(b-a)
	\]

	En este caso:

	\[
	\int_{-1}^{1}(1-x+x^2)\,dx
	=
	(1-c+c^2)(1-(-1))
	\]

	\[
	\left[
	x-\frac{x^2}{2}+\frac{x^3}{3}
	\right]_{-1}^{1}
	=
	2(1-c+c^2)
	\]

	\[
	\frac{8}{3}
	=
	2(1-c+c^2)
	\]

	\[
	\frac{4}{3}
	=
	1-c+c^2
	\]

	\[
	c^2-c-\frac{1}{3}=0
	\]

	Resolviendo la ecuación:

	\[
	c
	=
	\frac{1\pm\sqrt{1+\frac{4}{3}}}{2}
	=
	\frac{1\pm\sqrt{\frac{7}{3}}}{2}
	=
	\frac{3\pm\sqrt{21}}{6}
	\]

	Como debe cumplirse $c\in(-1,1)$, únicamente es válida:

	\[
	\boxed{
	c=\frac{3-\sqrt{21}}{6}
	}
	\]

\section{Ejercicio 10}
	Calcula

	\[
	\lim_{x\to0}
	\frac{\displaystyle\int_0^{x^2}\sin(\sqrt{t})\,dt}{x^3}
	\]

	\subsection{Solución}

	La expresión presenta una indeterminación $\frac{0}{0}$, por lo que aplicamos la regla de L'Hôpital.

	Por el Teorema Fundamental del Cálculo:

	\[
	\frac{d}{dx}
	\left(
	\int_0^{x^2}\sin(\sqrt{t})\,dt
	\right)
	=
	\sin(\sqrt{x^2})\cdot 2x
	=
	2x\sin(|x|)
	\]

	Por tanto:

	\[
	\lim_{x\to0}
	\frac{\displaystyle\int_0^{x^2}\sin(\sqrt{t})\,dt}{x^3}
	=
	\lim_{x\to0}
	\frac{2x\sin(|x|)}{3x^2}
	=
	\lim_{x\to0}
	\frac{2\sin(|x|)}{3x}
	\]

	Estudiamos los límites laterales:

	\[
	\lim_{x\to0^+}
	\frac{2\sin(|x|)}{3x}
	=
	\frac{2}{3}
	\]

	\[
	\lim_{x\to0^-}
	\frac{2\sin(|x|)}{3x}
	=
	-\frac{2}{3}
	\]

	Como los límites laterales son distintos:

	\[
	\boxed{
	\lim_{x\to0}
	\frac{\displaystyle\int_0^{x^2}\sin(\sqrt{t})\,dt}{x^3}
	\text{ no existe}
	}
	\]

\section{Ejercicio 11}
	Calcula

	\[
	\int_1^{+\infty}\frac{1}{x^2}\,dx
	\]

	\subsection{Solución}

	Al tratarse de una integral impropia de primera especie:

	\[
	\int_1^{+\infty}\frac{1}{x^2}\,dx
	=
	\lim_{k\to+\infty}
	\int_1^k\frac{1}{x^2}\,dx
	\]

	\[
	=
	\lim_{k\to+\infty}
	\left[
	-\frac{1}{x}
	\right]_1^k
	=
	\lim_{k\to+\infty}
	\left(
	-\frac{1}{k}+1
	\right)
	\]

	\[
	=
	1
	\]

	Por tanto, la integral es convergente y:

	\[
	\boxed{
	\int_1^{+\infty}\frac{1}{x^2}\,dx=1
	}
	\]

\section{Ejercicio 12}
	Dada la integral impropia

	\[
	\int_1^{+\infty}\frac{1}{x^\alpha}\,dx
	\]

	¿para qué valores $\alpha\in\mathbb{R}$ la integral es convergente?
	¿Para qué valores es divergente?

	\subsection{Solución}

	Expresamos la integral impropia mediante un límite:

	\[
	\int_1^{+\infty}\frac{1}{x^\alpha}\,dx
	=
	\lim_{k\to+\infty}
	\int_1^k x^{-\alpha}\,dx
	\]

	Para $\alpha\neq1$:

	\[
	=
	\lim_{k\to+\infty}
	\left[
	\frac{x^{1-\alpha}}{1-\alpha}
	\right]_1^k
	=
	\lim_{k\to+\infty}
	\frac{k^{1-\alpha}-1}{1-\alpha}
	\]

	Si $\alpha>1$:

	\[
	1-\alpha<0
	\qquad\Longrightarrow\qquad
	\lim_{k\to+\infty}k^{1-\alpha}=0
	\]

	\[
	\int_1^{+\infty}\frac{1}{x^\alpha}\,dx
	=
	\frac{1}{\alpha-1}
	\]

	Por tanto, la integral es convergente para:

	\[
	\boxed{\alpha>1}
	\]

	Si $\alpha<1$:

	\[
	1-\alpha>0
	\qquad\Longrightarrow\qquad
	\lim_{k\to+\infty}k^{1-\alpha}=+\infty
	\]

	por lo que la integral es divergente.

	Para $\alpha=1$:

	\[
	\int_1^{+\infty}\frac{1}{x}\,dx
	=
	\lim_{k\to+\infty}
	\left[\ln(x)\right]_1^k
	=
	\lim_{k\to+\infty}\ln(k)
	=
	+\infty
	\]

	Por tanto:

	\[
	\boxed{
	\begin{cases}
	\text{Convergente}, & \alpha>1 \\[2mm]
	\text{Divergente}, & \alpha\leq1
	\end{cases}
	}
	\]

\section{Ejercicio 13}
	Determina la convergencia o divergencia de

	\[
	\int_2^{+\infty}\frac{x+1}{x^2}\,dx
	\]

	\subsection{Solución}

	Descomponemos el integrando:

	\[
	\frac{x+1}{x^2}
	=
	\frac{1}{x}
	+
	\frac{1}{x^2}
	\]

	Por tanto:

	\[
	\int_2^{+\infty}\frac{x+1}{x^2}\,dx
	=
	\lim_{k\to+\infty}
	\int_2^k
	\left(
	\frac{1}{x}+\frac{1}{x^2}
	\right)\,dx
	\]

	\[
	=
	\lim_{k\to+\infty}
	\left[
	\ln(x)-\frac{1}{x}
	\right]_2^k
	\]

	\[
	=
	\lim_{k\to+\infty}
	\left(
	\ln(k)-\frac{1}{k}
	-\ln(2)+\frac{1}{2}
	\right)
	\]

	Como:

	\[
	\lim_{k\to+\infty}\ln(k)=+\infty
	\]

	se obtiene:

	\[
	\boxed{
	\int_2^{+\infty}\frac{x+1}{x^2}\,dx
	\text{ es divergente}
	}
	\]

\section{Ejercicio 14}
	Determina la convergencia o divergencia de

	\[
	\int_0^{+\infty}\cos(x)\,dx
	\]

	\subsection{Solución}

	Al tratarse de una integral impropia de primera especie:

	\[
	\int_0^{+\infty}\cos(x)\,dx
	=
	\lim_{k\to+\infty}
	\int_0^k\cos(x)\,dx
	\]

	\[
	=
	\lim_{k\to+\infty}
	\left[
	\sin(x)
	\right]_0^k
	=
	\lim_{k\to+\infty}\sin(k)
	\]

	El límite no existe, ya que $\sin(k)$ oscila entre $-1$ y $1$.

	Por tanto:

	\[
	\boxed{
	\int_0^{+\infty}\cos(x)\,dx
	\text{ es divergente}
	}
	\]

\section{Ejercicio 15}
	Determina la convergencia o divergencia de

	\[
	\int_0^1\frac{1}{\sqrt{x}}\,dx
	\]

	\subsection{Solución}

	La función presenta una discontinuidad infinita en $x=0$, por lo que se trata de una integral impropia de segunda especie:

	\[
	\int_0^1\frac{1}{\sqrt{x}}\,dx
	=
	\lim_{\varepsilon\to0^+}
	\int_{\varepsilon}^1 x^{-\frac{1}{2}}\,dx
	\]

	\[
	=
	\lim_{\varepsilon\to0^+}
	\left[
	2\sqrt{x}
	\right]_{\varepsilon}^1
	=
	\lim_{\varepsilon\to0^+}
	\left(
	2-2\sqrt{\varepsilon}
	\right)
	\]

	\[
	=
	2
	\]

	Por tanto, la integral es convergente y:

	\[
	\boxed{
	\int_0^1\frac{1}{\sqrt{x}}\,dx=2
	}
	\]

\section{Ejercicio 16}
	Calcula

	\[
	\int_0^1\frac{1}{x}\,dx
	\]

	\subsection{Solución}

	La función presenta una discontinuidad infinita en $x=0$, por lo que se trata de una integral impropia de segunda especie:

	\[
	\int_0^1\frac{1}{x}\,dx
	=
	\lim_{\varepsilon\to0^+}
	\int_{\varepsilon}^1\frac{1}{x}\,dx
	\]

	\[
	=
	\lim_{\varepsilon\to0^+}
	\left[
	\ln(x)
	\right]_{\varepsilon}^1
	=
	\lim_{\varepsilon\to0^+}
	\left(
	-\ln(\varepsilon)
	\right)
	\]

	\[
	=
	+\infty
	\]

	Por tanto:

	\[
	\boxed{
	\int_0^1\frac{1}{x}\,dx
	\text{ es divergente}
	}
	\]

\section{Ejercicio 17}
	Calcula

	\[
	\int_0^1\ln(x)\,dx
	\]

	\subsection{Solución}

	La función presenta una discontinuidad infinita en $x=0$, por lo que se trata de una integral impropia de segunda especie:

	\[
	\int_0^1\ln(x)\,dx
	=
	\lim_{\varepsilon\to0^+}
	\int_{\varepsilon}^1\ln(x)\,dx
	\]

	Integramos por partes:

	\[
	\int\ln(x)\,dx
	=
	x\ln(x)-x
	\]

	Por tanto:

	\[
	\int_0^1\ln(x)\,dx
	=
	\lim_{\varepsilon\to0^+}
	\left[
	x\ln(x)-x
	\right]_{\varepsilon}^1
	\]

	\[
	=
	\lim_{\varepsilon\to0^+}
	\left(
	-1-\varepsilon\ln(\varepsilon)+\varepsilon
	\right)
	\]

	Como:

	\[
	\lim_{\varepsilon\to0^+}
	\varepsilon\ln(\varepsilon)=0
	\]

	se obtiene:

	\[
	\boxed{
	\int_0^1\ln(x)\,dx=-1
	}
	\]

\section{Ejercicio 18}
	Calcula

	\[
	\int_0^3\frac{1}{(x-1)^{2/3}}\,dx
	\]

	y compara el resultado con su valor principal.

	\subsection{Solución}

	La función presenta una discontinuidad infinita en $x=1$, por lo que dividimos la integral:

	\[
	\int_0^3\frac{1}{(x-1)^{2/3}}\,dx
	=
	\lim_{\varepsilon_1\to0^+}
	\int_0^{1-\varepsilon_1}\frac{1}{(x-1)^{2/3}}\,dx
	+
	\lim_{\varepsilon_2\to0^+}
	\int_{1+\varepsilon_2}^3\frac{1}{(x-1)^{2/3}}\,dx
	\]

	Como:

	\[
	\int\frac{1}{(x-1)^{2/3}}\,dx
	=
	3(x-1)^{1/3}
	\]

	tenemos:

	\[
	\lim_{\varepsilon_1\to0^+}
	3\left[(x-1)^{1/3}\right]_0^{1-\varepsilon_1}
	=
	\lim_{\varepsilon_1\to0^+}
	3\left(1-\sqrt[3]{\varepsilon_1}\right)
	=
	3
	\]

	\[
	\lim_{\varepsilon_2\to0^+}
	3\left[(x-1)^{1/3}\right]_{1+\varepsilon_2}^3
	=
	\lim_{\varepsilon_2\to0^+}
	3\left(\sqrt[3]{2}-\sqrt[3]{\varepsilon_2}\right)
	=
	3\sqrt[3]{2}
	\]

	Por tanto:

	\[
	\boxed{
	\int_0^3\frac{1}{(x-1)^{2/3}}\,dx
	=
	3\left(1+\sqrt[3]{2}\right)
	}
	\]

	Calculamos ahora el valor principal:

	\[
	\operatorname{V.P.}
	\left(
	\int_0^3\frac{1}{(x-1)^{2/3}}\,dx
	\right)
	=
	\lim_{\varepsilon\to0^+}
	\left(
	\int_0^{1-\varepsilon}\frac{1}{(x-1)^{2/3}}\,dx
	+
	\int_{1+\varepsilon}^3\frac{1}{(x-1)^{2/3}}\,dx
	\right)
	\]

	\[
	=
	\lim_{\varepsilon\to0^+}
	3\left(
	1+\sqrt[3]{2}-2\sqrt[3]{\varepsilon}
	\right)
	=
	3\left(1+\sqrt[3]{2}\right)
	\]

	Por tanto, la integral impropia converge y coincide con su valor principal:

	\[
	\boxed{
	\operatorname{V.P.}
	=
	\int_0^3\frac{1}{(x-1)^{2/3}}\,dx
	=
	3\left(1+\sqrt[3]{2}\right)
	}
	\]

\section{Ejercicio 19}
	Calcula

	\[
	\int_{-1}^{1}\frac{1}{x^3}\,dx
	\]

	y compara el resultado con su valor principal.

	\subsection{Solución}

	La función presenta una discontinuidad infinita en $x=0$, por lo que dividimos la integral:

	\[
	\int_{-1}^{1}\frac{1}{x^3}\,dx
	=
	\lim_{\varepsilon_1\to0^+}
	\int_{-1}^{-\varepsilon_1}\frac{1}{x^3}\,dx
	+
	\lim_{\varepsilon_2\to0^+}
	\int_{\varepsilon_2}^{1}\frac{1}{x^3}\,dx
	\]

	Como:

	\[
	\int\frac{1}{x^3}\,dx
	=
	-\frac{1}{2x^2}
	\]

	tenemos:

	\[
	\lim_{\varepsilon_1\to0^+}
	\left[
	-\frac{1}{2x^2}
	\right]_{-1}^{-\varepsilon_1}
	=
	\lim_{\varepsilon_1\to0^+}
	\left(
	-\frac{1}{2\varepsilon_1^2}
	+\frac{1}{2}
	\right)
	=
	-\infty
	\]

	\[
	\lim_{\varepsilon_2\to0^+}
	\left[
	-\frac{1}{2x^2}
	\right]_{\varepsilon_2}^{1}
	=
	\lim_{\varepsilon_2\to0^+}
	\left(
	-\frac{1}{2}
	+\frac{1}{2\varepsilon_2^2}
	\right)
	=
	+\infty
	\]

	Por tanto, la integral impropia es divergente:

	\[
	\boxed{
	\int_{-1}^{1}\frac{1}{x^3}\,dx
	\text{ es divergente}
	}
	\]

	Calculamos ahora el valor principal:

	\[
	\operatorname{V.P.}
	\left(
	\int_{-1}^{1}\frac{1}{x^3}\,dx
	\right)
	=
	\lim_{\varepsilon\to0^+}
	\left(
	\int_{-1}^{-\varepsilon}\frac{1}{x^3}\,dx
	+
	\int_{\varepsilon}^{1}\frac{1}{x^3}\,dx
	\right)
	\]

	\[
	=
	\lim_{\varepsilon\to0^+}
	\left[
	-\frac{1}{2\varepsilon^2}
	+\frac{1}{2}
	-\frac{1}{2}
	+\frac{1}{2\varepsilon^2}
	\right]
	=
	0
	\]

	Por tanto:

	\[
	\boxed{
	\operatorname{V.P.}
	\left(
	\int_{-1}^{1}\frac{1}{x^3}\,dx
	\right)
	=
	0
	}
	\]

\section{Ejercicio 20}
	Calcula

	\[
	\int_{-2}^{4}\frac{1}{x^2}\,dx
	\]

	de forma directa y mediante el cambio $x=\frac{1}{t}$.

	\subsection{Solución}

	La función presenta una discontinuidad infinita en $x=0$, por lo que dividimos la integral:

	\[
	\int_{-2}^{4}\frac{1}{x^2}\,dx
	=
	\lim_{\varepsilon_1\to0^+}
	\int_{-2}^{-\varepsilon_1}\frac{1}{x^2}\,dx
	+
	\lim_{\varepsilon_2\to0^+}
	\int_{\varepsilon_2}^{4}\frac{1}{x^2}\,dx
	\]

	\[
	=
	\lim_{\varepsilon_1\to0^+}
	\left[
	-\frac{1}{x}
	\right]_{-2}^{-\varepsilon_1}
	+
	\lim_{\varepsilon_2\to0^+}
	\left[
	-\frac{1}{x}
	\right]_{\varepsilon_2}^{4}
	\]

	\[
	=
	\lim_{\varepsilon_1\to0^+}
	\left(
	\frac{1}{\varepsilon_1}-\frac{1}{2}
	\right)
	+
	\lim_{\varepsilon_2\to0^+}
	\left(
	\frac{1}{\varepsilon_2}-\frac{1}{4}
	\right)
	=
	+\infty
	\]

	Por tanto:

	\[
	\boxed{
	\int_{-2}^{4}\frac{1}{x^2}\,dx
	\text{ es divergente}
	}
	\]

	Realizamos ahora el cambio:

	\[
	x=\frac{1}{t}
	\qquad\Longrightarrow\qquad
	dx=-\frac{1}{t^2}\,dt
	\]

	\[
	\frac{1}{x^2}\,dx
	=
	t^2\left(-\frac{1}{t^2}\right)\,dt
	=
-dt
	\]

	Como el cambio no está definido en $x=0$, tratamos ambos intervalos por separado:

	\[
	\int_{-2}^{0}\frac{1}{x^2}\,dx
	=
	\int_{-\infty}^{-\frac{1}{2}}dt
	=
	+\infty
	\]

	\[
	\int_{0}^{4}\frac{1}{x^2}\,dx
	=
	\int_{\frac{1}{4}}^{+\infty}dt
	=
	+\infty
	\]

	Por tanto, mediante el cambio de variable obtenemos igualmente:

	\[
	\boxed{
	\int_{-2}^{4}\frac{1}{x^2}\,dx
	\text{ es divergente}
	}
	\]

\section{Ejercicio 21}
	Calcula

	\[
	\int_0^1\sqrt{\frac{1-x}{x}}\,dx
	\]

	\subsection{Solución}

	Expresamos la integral en la forma de la función Beta:

	\[
	\int_0^1\sqrt{\frac{1-x}{x}}\,dx
	=
	\int_0^1
	x^{-\frac{1}{2}}
	(1-x)^{\frac{1}{2}}
	\,dx
	\]

	Recordando que:

	\[
	\beta(p,q)
	=
	\int_0^1
	x^{p-1}(1-x)^{q-1}\,dx
	\]

	identificamos:

	\[
	p-\frac{}{}1=\!-\frac{1}{2}
	\qquad\Longrightarrow\qquad
	p=\frac{1}{2}
	\]

	\[
	q-1=\frac{1}{2}
	\qquad\Longrightarrow\qquad
	q=\frac{3}{2}
	\]

	Por tanto:

	\[
	\int_0^1\sqrt{\frac{1-x}{x}}\,dx
	=
	\beta\left(\frac{1}{2},\frac{3}{2}\right)
	\]

	Utilizamos la relación:

	\[
	\beta(p,q)
	=
	\frac{\Gamma(p)\Gamma(q)}
	{\Gamma(p+q)}
	\]

	\[
	=
	\frac{
	\Gamma\left(\frac{1}{2}\right)
	\Gamma\left(\frac{3}{2}\right)
	}{
	\Gamma(2)
	}
	\]

	Como:

	\[
	\Gamma\left(\frac{1}{2}\right)=\sqrt{\pi},
	\qquad
	\Gamma\left(\frac{3}{2}\right)
	=
	\frac{1}{2}\Gamma\left(\frac{1}{2}\right)
	=
	\frac{\sqrt{\pi}}{2},
	\qquad
	\Gamma(2)=1
	\]

	obtenemos:

	\[
	\boxed{
	\int_0^1\sqrt{\frac{1-x}{x}}\,dx
	=
	\frac{\pi}{2}
	}
	\]

\section{Ejercicio 22}
	Calcula

	\[
	\int_0^1\sqrt{1-x^5}\,dx
	\]

	\subsection{Solución}

	Realizamos el cambio:

	\[
	t=x^5
	\qquad\Longrightarrow\qquad
	x=t^{1/5},
	\qquad
	dx=\frac{1}{5}t^{-4/5}\,dt
	\]

	Entonces:

	\[
	\int_0^1\sqrt{1-x^5}\,dx
	=
	\frac{1}{5}
	\int_0^1
	t^{-4/5}(1-t)^{1/2}\,dt
	\]

	Recordando que:

	\[
	\beta(p,q)
	=
	\int_0^1
	t^{p-1}(1-t)^{q-1}\,dt
	\]

	identificamos:

	\[
	p=\frac{1}{5},
	\qquad
	q=\frac{3}{2}
	\]

	Por tanto:

	\[
	\int_0^1\sqrt{1-x^5}\,dx
	=
	\frac{1}{5}
	\beta\left(\frac{1}{5},\frac{3}{2}\right)
	\]

	Utilizando:

	\[
	\beta(p,q)
	=
	\frac{\Gamma(p)\Gamma(q)}
	{\Gamma(p+q)}
	\]

	\[
	=
	\frac{1}{5}
	\frac{
	\Gamma\left(\frac{1}{5}\right)
	\Gamma\left(\frac{3}{2}\right)
	}{
	\Gamma\left(\frac{17}{10}\right)
	}
	\]

	Como:

	\[
	\Gamma\left(\frac{3}{2}\right)
	=
	\frac{\sqrt{\pi}}{2}
	\]

	obtenemos:

	\[
	\boxed{
	\int_0^1\sqrt{1-x^5}\,dx
	=
	\frac{\sqrt{\pi}}{10}
	\frac{
	\Gamma\left(\frac{1}{5}\right)
	}{
	\Gamma\left(\frac{17}{10}\right)
	}
	}
	\]

\section{Ejercicio 23}
	Calcula

	\[
	\int_0^{+\infty}x^{4/3}e^{-x}\,dx
	\]

	\subsection{Solución}

	Recordando la definición de la función Gamma:

	\[
	\Gamma(p)
	=
	\int_0^{+\infty}x^{p-1}e^{-x}\,dx
	\]

	identificamos:

	\[
	p-1=\frac{4}{3}
	\qquad\Longrightarrow\qquad
	p=\frac{7}{3}
	\]

	Por tanto:

	\[
	\int_0^{+\infty}x^{4/3}e^{-x}\,dx
	=
	\Gamma\left(\frac{7}{3}\right)
	\]

	Utilizando la propiedad:

	\[
	\Gamma(p)
	=
	(p-1)\Gamma(p-1)
	\]

	\[
	\Gamma\left(\frac{7}{3}\right)
	=
	\frac{4}{3}
	\Gamma\left(\frac{4}{3}\right)
	=
	\frac{4}{3}\cdot\frac{1}{3}
	\Gamma\left(\frac{1}{3}\right)
	\]

	Por tanto:

	\[
	\boxed{
	\int_0^{+\infty}x^{4/3}e^{-x}\,dx
	=
	\frac{4}{9}
	\Gamma\left(\frac{1}{3}\right)
	}
	\]

\section{Ejercicio 24}
	Calcula mediante un cambio de variable la integral

	\[
	\int_1^{+\infty}\frac{\ln(x)}{x^2}\,dx
	\]

	\subsection{Solución}

	Realizamos el cambio:

	\[
	t=\frac{1}{x}
	\qquad\Longrightarrow\qquad
	x=\frac{1}{t},
	\qquad
	dx=-\frac{1}{t^2}\,dt
	\]

	Además:

	\[
	\ln(x)=\ln\left(\frac{1}{t}\right)=-\ln(t),
	\qquad
	\frac{1}{x^2}=t^2
	\]

	Los límites se transforman en:

	\[
	x=1\Longrightarrow t=1,
	\qquad
	x\to+\infty\Longrightarrow t\to0^+
	\]

	Por tanto:

	\[
	\int_1^{+\infty}\frac{\ln(x)}{x^2}\,dx
	=
	\int_1^0
	(-\ln(t))t^2
	\left(-\frac{1}{t^2}\right)\,dt
	=
	\int_1^0\ln(t)\,dt
	\]

	\[
	=
	-\int_0^1\ln(t)\,dt
	\]

	Como:

	\[
	\int_0^1\ln(t)\,dt=-1
	\]

	obtenemos:

	\[
	\boxed{
	\int_1^{+\infty}\frac{\ln(x)}{x^2}\,dx=1
	}
	\]

\section{Ejercicio 25}
	Calcula $\Gamma\left(\frac{1}{2}\right)$ utilizando para ello la función beta correspondiente.

	\subsection{Solución}

	Utilizamos la relación:

	\[
	\beta(p,q)
	=
	2\int_0^{\pi/2}
	(\cos t)^{2p-1}
	(\sin t)^{2q-1}\,dt
	\]

	Tomando:

	\[
	p=q=\frac{1}{2}
	\]

	obtenemos:

	\[
	\beta\left(\frac{1}{2},\frac{1}{2}\right)
	=
	2\int_0^{\pi/2}1\,dt
	=
	\pi
	\]

	Por otra parte:

	\[
	\beta(p,q)
	=
	\frac{\Gamma(p)\Gamma(q)}
	{\Gamma(p+q)}
	\]

	\[
	\beta\left(\frac{1}{2},\frac{1}{2}\right)
	=
	\frac{
	\Gamma\left(\frac{1}{2}\right)^2
	}{
	\Gamma(1)
	}
	=
	\Gamma\left(\frac{1}{2}\right)^2
	\]

	Por tanto:

	\[
	\Gamma\left(\frac{1}{2}\right)^2
	=
	\pi
	\]

	Como $\Gamma\left(\frac{1}{2}\right)>0$:

	\[
	\boxed{
	\Gamma\left(\frac{1}{2}\right)=\sqrt{\pi}
	}
	\]

\section{Ejercicio 26}
	Calcula la integral paramétrica

	\[
	\int_1^3(3x-1)\cos(tx)\,dx
	\]

	\subsection{Solución}

	Integramos por partes:

	\[
	u=3x-1,
	\qquad
	dv=\cos(tx)\,dx
	\]

	\[
	du=3\,dx,
	\qquad
	v=\frac{\sin(tx)}{t}
	\]

	Por tanto:

	\[
	\int(3x-1)\cos(tx)\,dx
	=
	\frac{(3x-1)\sin(tx)}{t}
	-
	\frac{3}{t}\int\sin(tx)\,dx
	\]

	\[
	=
	\frac{(3x-1)\sin(tx)}{t}
	+
	\frac{3\cos(tx)}{t^2}
	\]

	Aplicando los límites:

	\[
	\int_1^3(3x-1)\cos(tx)\,dx
	=
	\left[
	\frac{(3x-1)\sin(tx)}{t}
	+
	\frac{3\cos(tx)}{t^2}
	\right]_1^3
	\]

	\[
	=
	\frac{8\sin(3t)-2\sin(t)}{t}
	+
	\frac{3\left(\cos(3t)-\cos(t)\right)}{t^2},
	\qquad t\neq0
	\]

	Para $t=0$:

	\[
	\int_1^3(3x-1)\,dx
	=
	\left[
	\frac{3x^2}{2}-x
	\right]_1^3
	=
	10
	\]

	Por tanto:

	\[
	\boxed{
	F(t)=
	\begin{cases}
	\displaystyle
	\frac{8\sin(3t)-2\sin(t)}{t}
	+
	\frac{3\left(\cos(3t)-\cos(t)\right)}{t^2},
	& t\neq0 \\[4mm]
	10,
	& t=0
	\end{cases}
	}
	\]

\section{Ejercicio 27}
	Calcula

	\[
	\lim_{t\to0}
	\int_1^3(3x-1)\cos(tx)\,dx
	\]

	\subsection{Solución}

	Como el integrando es continuo, podemos intercambiar el límite y la integral:

	\[
	\lim_{t\to0}
	\int_1^3(3x-1)\cos(tx)\,dx
	=
	\int_1^3
	\lim_{t\to0}
	(3x-1)\cos(tx)\,dx
	\]

	\[
	=
	\int_1^3(3x-1)\,dx
	\]

	\[
	=
	\left[
	\frac{3x^2}{2}-x
	\right]_1^3
	=
	\left(\frac{27}{2}-3\right)
	-
	\left(\frac{3}{2}-1\right)
	\]

	\[
	=
	10
	\]

	\[
	\boxed{10}
	\]

\section{Ejercicio 28}
	Dada

	\[
	I(t)=\int_{t^2}^{t}\sin(x^2+t^2)\,dx
	\]

	determina la expresión de $I'(t)$.

	\subsection{Solución}

	Utilizamos la regla de derivación para integrales paramétricas:

	\[
	I'(t)
	=
	\int_{\alpha(t)}^{\beta(t)}
	\frac{\partial f(x,t)}{\partial t}\,dx
	+
	f(\beta(t),t)\beta'(t)
	-
	f(\alpha(t),t)\alpha'(t)
	\]

	En este caso:

	\[
	f(x,t)=\sin(x^2+t^2),
	\qquad
	\alpha(t)=t^2,
	\qquad
	\beta(t)=t
	\]

	\[
	\frac{\partial f(x,t)}{\partial t}
	=
	2t\cos(x^2+t^2)
	\]

	Por tanto:

	\[
	I'(t)
	=
	\int_{t^2}^{t}2t\cos(x^2+t^2)\,dx
	+
	\sin(t^2+t^2)
	-
	2t\sin(t^4+t^2)
	\]

	\[
	\boxed{
	I'(t)
	=
	2t\int_{t^2}^{t}\cos(x^2+t^2)\,dx
	+
	\sin(2t^2)
	-
	2t\sin(t^4+t^2)
	}
	\]

\section{Ejercicio 29}
	Utilizando derivación paramétrica, calcula

	\[
	\int_0^{\pi/2}
	\frac{\arctan(t\sin x)}{\sin x}\,dx
	\]

	\subsection{Solución}

	Definimos:

	\[
	I(t)
	=
	\int_0^{\pi/2}
	\frac{\arctan(t\sin x)}{\sin x}\,dx
	\]

	Derivamos respecto de $t$:

	\[
	I'(t)
	=
	\int_0^{\pi/2}
	\frac{1}{1+t^2\sin^2 x}\,dx
	\]

	Realizamos el cambio:

	\[
	u=\tan x,
	\qquad
	dx=\frac{du}{1+u^2},
	\qquad
	\sin^2 x=\frac{u^2}{1+u^2}
	\]

	\[
	I'(t)
	=
	\int_0^{+\infty}
	\frac{1}{1+(1+t^2)u^2}\,du
	\]

	\[
	=
	\left[
	\frac{1}{\sqrt{1+t^2}}
	\arctan\left(u\sqrt{1+t^2}\right)
	\right]_0^{+\infty}
	=
	\frac{\pi}{2\sqrt{1+t^2}}
	\]

	Integramos respecto de $t$:

	\[
	I(t)
	=
	\frac{\pi}{2}
	\int\frac{1}{\sqrt{1+t^2}}\,dt
	=
	\frac{\pi}{2}
	\ln\left(t+\sqrt{1+t^2}\right)+C
	\]

	Como:

	\[
	I(0)=0
	\]

	se obtiene:

	\[
	C=0
	\]

	Por tanto:

	\[
	\boxed{
	\int_0^{\pi/2}
	\frac{\arctan(t\sin x)}{\sin x}\,dx
	=
	\frac{\pi}{2}
	\ln\left(t+\sqrt{1+t^2}\right)
	}
	\]

\section{Ejercicio 30}
	Calcula por derivación paramétrica

	\[
	I(t)=\int_0^t\frac{1}{(x^2+t^2)^3}\,dx,
	\qquad t>0
	\]

	\subsection{Solución}

	Derivamos respecto de $t$:

	\[
	I'(t)
	=
	\frac{1}{(t^2+t^2)^3}
	+
	\int_0^t
	\frac{-6t}{(x^2+t^2)^4}\,dx
	\]

	\[
	I'(t)
	=
	\frac{1}{8t^6}
	-
	6t\int_0^t\frac{1}{(x^2+t^2)^4}\,dx
	\]

	Por otra parte:

	\[
	\frac{d}{dx}
	\left(
	\frac{x}{(x^2+t^2)^3}
	\right)
	=
	\frac{1}{(x^2+t^2)^3}
	-
	\frac{6x^2}{(x^2+t^2)^4}
	\]

	Integrando entre $0$ y $t$:

	\[
	\frac{1}{8t^5}
	=
	I(t)
	-
	6\int_0^t
	\frac{x^2}{(x^2+t^2)^4}\,dx
	\]

	Como:

	\[
	x^2=(x^2+t^2)-t^2
	\]

	se obtiene:

	\[
	\frac{1}{8t^5}
	=
	-5I(t)
	+
	6t^2\int_0^t
	\frac{1}{(x^2+t^2)^4}\,dx
	\]

	Por tanto:

	\[
	6t\int_0^t
	\frac{1}{(x^2+t^2)^4}\,dx
	=
	\frac{1}{8t^6}
	+
	\frac{5I(t)}{t}
	\]

	Sustituyendo en $I'(t)$:

	\[
	I'(t)
	=
	-\frac{5}{t}I(t)
	\]

	\[
	\frac{I'(t)}{I(t)}
	=
	-\frac{5}{t}
	\]

	\[
	\ln I(t)
	=
	-5\ln t+C
	\]

	\[
	I(t)=\frac{C}{t^5}
	\]

	Calculamos la constante tomando $t=1$:

	\[
	C=I(1)
	=
	\int_0^1\frac{1}{(1+x^2)^3}\,dx
	\]

	Con el cambio $x=\tan u$:

	\[
	C
	=
	\int_0^{\pi/4}\cos^4u\,du
	=
	\frac{3\pi}{32}+\frac{1}{4}
	=
	\frac{3\pi+8}{32}
	\]

	Por tanto:

	\[
	\boxed{
	I(t)
	=
	\frac{3\pi+8}{32t^5}
	}
	\]
\end{document}